% !TeX encoding = utf8
% !TeX TXS-program:bibliography = txs:///biber
%%%%%%%%%%%%%%
% Fichero: MemoriaTFG
% Autor: Jesús Salido Tercero (http://www.uclm.es/profesorado/jsalido)
% Fecha (creación): Febrero 2010 
% Rev. : Febrero 2017
% Descripción: Plantilla para memoria de TFG 
% (Escuela Sup. de Informática, UCLM). Creada para el curso 
% “LaTeX esencial para preparación de TFG, Tesis y otros documentos 
% académicos” (Esc. Sup. Informática-UCLM)
%
% Comentarios: Preparada para `pdflatex' y `biblatex' (con `biber'). 
% Documento editado con TeXstudio. 
% Para su compilación se aconseja utilizar como compilador 
% por defecto `latexmk.'
%%%%%%%%%%%%%%

% EDITAR: Idioma pral. \spanishtrue (Español), \spanishfalse (Inglés)
\newif\ifspanish
\spanishtrue
%\spanishfalse

\documentclass[11pt,a4paper,twoside,openright,final]{book}
\usepackage[utf8]{inputenx} % Codificación de entrada
\usepackage[spanish,english]{babel} % Internacionalización
\usepackage{indentfirst} % Para asegurar sangrado en 1ª línea tras sección (necesario con varios idiomas)
\usepackage[top=2.5cm,bottom=2.5cm,inner=3.5cm,outer=2cm]{geometry}


% Tipografía
\usepackage{libertine}
\usepackage[libertine]{newtxmath}

% Inclusión de símbolos
\usepackage{textcomp,marvosym,pifont} % Símbolos
\usepackage{cclicenses} % Para símbolos de licencias CC
\usepackage{amsmath,amsfonts,amssymb} % Caracteres matemáticos
\usepackage[T1]{fontenc} % Codificación de salida 
\usepackage{microtype} % Mejoras tipográficas para pdflatex

\usepackage{url} % Para escritura de URL
\urlstyle{sf}
\usepackage[bookmarks,bookmarksnumbered=true,hyperfootnotes=false,hidelinks]{hyperref}

% Definición de colores
% OJO: Este paquete debe cargarse antes de ctable. 
\usepackage[usenames,dvipsnames,svgnames,x11names,table]{xcolor}
\definecolor{sombra}{gray}{.70}


% Tablas y gráficos
\usepackage{framed} % Recuadro en torno a texto 
\usepackage{fancybox} % Recuadro exóticos
\usepackage{graphicx}  % Inclusión de figuras
\usepackage{subfigure} % Inclusión de subfiguras
\graphicspath{{./figs/}} % Path de búsqueda de ficheros gráficos
\DeclareGraphicsExtensions{.pdf,.png,.jpg} % Precedencia de extensiones
\usepackage{rotating}  % Giro de cajas (texto, figuras, tablas) (No DVI)
\usepackage{tikz}
\usepackage{ctable} % Inclusión de tablas. (ctable incluye): color,xkeyval,array,tabularx,booktabs,rotating
\usepackage{multirow}


% Paquetes especiales para Informática
\usepackage{listings} % Inclusión de listados de código
% Algoritmos
\ifspanish
\usepackage[lined,boxruled,algochapter,commentsnumbered,spanish]{algorithm2e}
\else
\usepackage[lined,boxruled,algochapter,commentsnumbered,english]{algorithm2e}
\fi
% Paquete específico para la generación de algorithmos en LaTeX. Por su complejidad se requiere consultar la documentación del paquete para su correcto uso.


% Ajustes de formato
\usepackage{paralist, multicol} % Mayor control de listas
\usepackage{titlesec} % Personalización completa de títulos de secciones
\usepackage{sectsty} % Personalización de títulos de sección con una interfaz más simple que la suministrada por titlesec
% Tipo de letra empleado en títulos de secciones (paquete sectsty)
\usepackage[margin=10pt,font=small,labelfont=bf,format=hang]{caption} % Personalización de títulos de figuras y tablas
\usepackage{fancyhdr} % Cabeceras y pies personalizados


% Bibliografía (Multilingüe en español e inglés, empleando el idioma de la fuente)
\usepackage[backend=biber,sortcites,autolang=other,language=auto]{biblatex}
% Línea añadida para eliminar el idioma de la fuente bibliográfica.
\AtEveryBibitem{\clearfield{note} \clearlist{language}}
\addbibresource{biblioTFG.bib}
\usepackage[autostyle]{csquotes}

 % Paquete para índice alfabético
\usepackage{makeidx}
\makeindex

% Comandos y ajustes de usuario
% Ajustes definidos por el usuario

% Control de espaciado entre párrafos e indentación
\setlength{\parskip}{1.3ex plus 0.2ex minus 0.2ex} % Espacio entre párrafos (más versatil cuando se fija como una rubber length (longitud elástica)
%\setlength{\parskip}{2mm plus 0.2mm minus 0.2mm} (Clase esi-tfg.cls por D. Villa)
%\setlength{\parindent}{8ex} % Valor por defecto de la indentación en español. 
%Sólo es preciso activar para modificar el valor por defecto (indicar valor cero para anular)


% Ajuste de cabecera y pie de página (paquete fancyhdr)
\fancyhf{} % Reset de la cabecera y pie
% En las páginas impares, parte izquierda del encabezado, aparecerá el nombre de capítulo
\fancyhead[LO]{\sffamily\leftmark} 
% En las páginas pares, parte derecha del encabezado, aparecerá el nombre de sección
\fancyhead[RE]{\sffamily\rightmark} 
% Números de página en las esquinas de los encabezados
\fancyhead[RO,LE]{\sffamily\thepage}

% Formato para el capítulo: N. Nombre
\renewcommand{\chaptermark}[1]{\markboth{\textbf{\thechapter.#1}}{}}
% Formato para la sección: N.M. Nombre
\renewcommand{\sectionmark}[1]{\markright{\textbf{\thesection. #1}}} 

% Ancho de la línea horizontal bajo el encabezado
\renewcommand{\headrulewidth}{0.6pt} 
% Ancho de la línea horizontal sobre el pie (en este ejemplo está vacío)
%\renewcommand{\footrulewidth}{0.6pt} 

%\setlength{\headheight}{16pt} % Por defecto \headheight 12pt, pero se agranda al emplear fancyhdr
\setlength{\headheight}{1.5\headheight} % Aumenta la altura del encabezado en una vez y media


% Evita que la última página de cap. tenga cabecera y pie 
% si dicha página está en blanco (para clase book)
\makeatletter
\def\cleardoublepage{\clearpage\if@twoside
\ifodd\c@page
\else\hbox{}\thispagestyle{empty}\newpage
\if@twocolumn\hbox{}\newpage\fi\fi\fi}
\makeatother



%====================================
% Configuración inicio capítulos (paquete titlesec)
\newcommand{\bigrule}{\titlerule[0.5mm]}

\titleformat{\chapter}[display] % cambiamos el formato de los capítulos
{\bfseries\Huge\sffamily} % por defecto se usarán caracteres de tamaño \Huge en negrita y familia sans serif
{% contenido de la etiqueta
\titlerule % línea horizontal
\filleft % texto alineado a la derecha
\Large{\MakeUppercase\chaptertitlename}\ % "Capítulo" o "Apéndice" en tamaño \Large en lugar de \Huge
\Large\thechapter} % número de capítulo en tamaño \Large
{0mm} % espacio mínimo entre etiqueta y cuerpo
{\filleft\MakeUppercase} % texto del cuerpo alineado a la derecha
[\vspace{0.5mm} \bigrule] % después del cuerpo, dejar espacio vertical y trazar línea horizontal gruesa

%====================================
% Tipo de letra empleado en títulos de secciones (paquete sectsty)
\sectionfont{\sffamily\bfseries\MakeUppercase}
\subsectionfont{\sffamily\bfseries}
\subsubsectionfont{\sffamily\bfseries}

%====================================
% Alguna definiciones de nivel de gris
\definecolor{gris97}{gray}{.97}
\definecolor{gris95}{gray}{.95}
\definecolor{gris75}{gray}{.75}
\definecolor{gris45}{gray}{.45}


% Personalización del entorno lstlisting (ver documentación paquete listings para más infor.)
   
\lstset{ % Estilo por defecto
	basicstyle={\footnotesize\ttfamily}, % Estilo básico para el texto
	%stringstyle=\textsl,        % Estilo para las cadenas
	stringstyle={\color{Red1}\ttfamily\bfseries},
	commentstyle={\color{Green4}\sffamily\bfseries},% Estilo para los comentarios
	keywordstyle={\color{Blue1}\bfseries},% Estilo para las palabras clave
	%	keywordstyle=[1]\textbf,    % Posibilidad de particularizar el estilo 
	%	keywordstyle=[2]\textbf,    %
	%	keywordstyle=[3]\textbf,    %
	%	keywordstyle=[4]\textbf,    %
	%	deletekeywords={}, 			% Quita keywords separadas por comas
	captionpos=t,               % Ajusta la posición de títulos 
	numbers=left,               % Posición de números de línea
	numberstyle={\tiny\sffamily\bfseries},          % Tamaño del número de línea
	numberfirstline=false,
	firstnumber=1, 				%  Nº de la primera línea
	stepnumber=1,               % Paso de línea numerada
	numbersep=10pt,             % Separación del texto al número de línea
	tabsize=2,                  % Tamaño del tabulador
	extendedchars=\true,        % Gestiona en empleo de caracteres extendidos utf8
	texcl=true,				    % Necesario para unicode en los comentarios
	breaklines=true,            % Ajusta división automática de líneas
	breakatwhitespace=true,     %
	frame=single,               % none, leftline, topline, bottomline, lines, single, shadowbox 
	frameround=tttt, 			% Redondea los bordes del frame
	rulecolor={\color{IndianRed}},    % Color del frame
	showspaces=false,           % Muestra espacios en blanco
	showtabs=false,             % Muestra tabuladores
	showstringspaces=true,      % Muestra espacios en blanco en las cadenas        
	xleftmargin=1cm,xrightmargin=1cm,
	breaklines=true,
	%	framexleftmargin=17pt
	%	framexrightmargin=5pt,
	%	framexbottommargin=4pt,
	backgroundcolor={\color{Azure1}} % Color del fondo
}


%Estilo definido para lenguaje C
\lstdefinestyle{C}{%
	language=C,
	frame=L,
	rulesep=.1pt,
	rulecolor=\color{black},
}


% Estilo definido para comandos de consola
\lstdefinestyle{consola}{%
	basicstyle={\color{White}\scriptsize\bf\ttfamily},
	backgroundcolor={\color{Black}},
	frame=none,
	showspaces=true
}




% Comandos definidos por el usuario
%====================================
% Comandos que definen datos del documento
\makeatletter
\newcommand{\tituloPrimera}[1]{\newcommand{\@tituloPrimera}{#1}}
\newcommand{\tituloSegunda}[1]{\newcommand{\@tituloSegunda}{#1}}
\newcommand{\titulo}[1]{\newcommand{\@titulo}{#1}\renewcommand{\@title}{#1}}
\newcommand{\tipoDoc}[1]{\newcommand{\@tipoDoc}{#1}}
\newcommand{\autor}[1]{\newcommand{\@autor}{#1}\renewcommand{\@author}{#1}}
\newcommand{\email}[1]{\newcommand{\@email}{\url{#1}}}
\newcommand{\director}[1]{\newcommand{\@director}{#1}}
\newcommand{\codirector}[1]{\newcommand{\@codirector}{#1}}
\newcommand{\tutor}[1]{\newcommand{\@tutor}{#1}}
\newcommand{\instEdu}[1]{\newcommand{\@instEdu}{#1}}
\newcommand{\centroEdu}[1]{\newcommand{\@centroEdu}{#1}}
\newcommand{\deptoEduPrimera}[1]{\newcommand{\@deptoEduPrimera}{#1}}
\newcommand{\deptoEduSegunda}[1]{\newcommand{\@deptoEduSegunda}{#1}}
\newcommand{\escudo}[1]{\newcommand{\@escudo}{#1}}
\newcommand{\titulacion}[1]{\newcommand{\@titulacion}{#1}}
\newcommand{\especialidad}[1]{\newcommand{\@especialidad}{#1}}
\newcommand{\fechaDef}[1]{\newcommand{\@fechaDef}{#1}}
\newcommand{\mesDef}[1]{\newcommand{\@mesDef}{#1}}
\newcommand{\yearDef}[1]{\newcommand{\@yearDef}{#1}}
\newcommand{\lugarDef}[1]{\newcommand{\@lugarDef}{#1}}
\makeatother


%====================================
% Portada (tfg)
% Portada (aquí gralmente no habrá que editar nada)
\makeatletter
\newcommand{\portadaTFG}{%
	\begin{titlepage}
		\begin{center}
			\includegraphics[width=3.5cm]{\@escudo}\vspace{1cm}
			
			{\LARGE \textbf{\@instEdu\\[1.5\parskip]
					\@centroEdu\\[2cm]
					\@titulacion}}\\[0.5cm]
			{\large \textbf{\@especialidad}}\\[1.5cm]
			{\LARGE \textbf{\@tipoDoc}}\\[1cm]
			
			{\LARGE \@tituloPrimera}\\ \smallskip%			
			\ifdefined\@tituloSegunda{\LARGE	\@tituloSegunda}\\[3cm]
			\else \vspace{3cm}
			\fi
			
			{\Large \@autor}\vfill%
		\end{center}
		
		\begin{flushright}
			{\Large \@fechaDef}
		\end{flushright}
		
		\cleardoublepage
\end{titlepage}}
\makeatother

%====================================
% Página inicial (tfg)
% Página inicial (es como la portada, añade Director/es o Tutor)
\makeatletter
\newcommand{\portadillaTFG}{%
	\begin{center}
		\includegraphics[width=3.5cm]{\@escudo}\\[1.5cm]
		
		{\LARGE \textbf{\@instEdu \\[1.5\parskip]
				\@centroEdu}}\\[0.5cm]
		{\Large \textbf{\@deptoEduPrimera}}\\ \smallskip%
		\ifdefined\@deptoEduSegunda{\Large \textbf{\@deptoEduSegunda}}\\[0.5cm]
		\else \vspace{0.5cm}
		\fi
		{\large \textbf{\@especialidad}}\\[1.5cm]
		
		{\LARGE \textbf{\@tipoDoc}}\\[1cm]
		
		
		{\LARGE \textbf{\@tituloPrimera}}\\ \smallskip%		
		\ifdefined\@tituloSegunda{\LARGE \textbf{\@tituloSegunda}}\\[3cm]
		\else \vspace{3cm}
		\fi
	\end{center}
	
	\begin{flushleft}
		{\Large Autor(a): \@autor} \\ \bigskip%
		{\Large Director(a): \@director} \\ \bigskip%
		% Si hay definido un codirector se añade automáticamente la línea siguiente
		\ifdefined\@codirector {\Large Director(a): \@codirector} \fi 
	\end{flushleft}
	\vfill%
	
	\begin{flushright}
		{\Large \@fechaDef}
	\end{flushright}
	
	\newpage}
\makeatother

%====================================
% Página tribunal (tfg)
% Tribunal
\makeatletter
\newcommand{\tribunalTFG}{
	{\flushright \LARGE \textsc{Tribunal:}}
	
	\vspace*{\stretch{0.5}}
	\hspace*{1cm}{\Large Presidente: \hrulefill}
	
	\vspace*{\stretch{0.5}}
	\hspace*{1cm}{\Large Vocal: \hrulefill}
	
	\vspace*{\stretch{0.5}}
	\hspace*{1cm}{\Large Secretario: \hrulefill}
	
	\vspace*{\stretch{0.5}}
	{\flushright \LARGE \textsc{Fecha de defensa:} \hrulefill}
	
	\vspace*{\stretch{1.5}}
	{\flushright \LARGE \textsc{Calificación:} \hrulefill}
	
	\vspace*{\stretch{2.5}}
	\begin{center}
		\begin{tabularx}{\linewidth}{X X X}
			{\large \textsc{Presidente}} & {\large \textsc{Vocal}} & {\large \textsc{Secretario}}\\[2.5cm]
			Fdo.: & Fdo.: & Fdo.:		
		\end{tabularx}
	\end{center}
	\cleardoublepage}
\makeatother



%====================================
% Comando para incluir la dedicatoria
% Con la opción stretch se puede colocar verticalmente la dedicatoria de forma relativa para que quede el doble de espacio por debajo que por encima
\newcommand{\dedicado}[1]{ % Dedicatoria
	\null\vspace{\stretch{1}}
	\begin{flushright}
	\emph{#1}
	\end{flushright}
	\vspace{\stretch{2}}\null
	\cleardoublepage
}



%====================================
% Créditos y licencia (opcional, 1 pág.)
% Todas las obras en gral. deberían presentar información relativa a la propiedad intelectual del contenido y condiciones bajo las cuales se puede distribuir y reproducir
\makeatletter
\newcommand{\creditos}[2]{%
	\null\vspace{6cm}
	{\small \noindent \@titulo\\
	\textcopyright{} \@autor, \@yearDef\\[1cm]
	#1}\\
	\includegraphics[width=0.15\linewidth]{#2}
	
	\clearpage
}
\makeatother




%====================================
% Añade el comando \tecla para crear indicaciones de pulsación de teclas
\usetikzlibrary{shadows} % Necesario para poder crear nuevo comando de indicación de pulsación de tecla.
\newcommand*\tecla[1]{%   
  \tikz[baseline=(key.base)]
    \node[%
      draw,
      fill=white,
      drop shadow={%
	      shadow xshift=0.25ex,
	      shadow yshift=-0.25ex,
	      fill=black,
	      opacity=0.75
      },
      rectangle,
      rounded corners=2pt,
      inner sep=1pt,
      line width=0.5pt,
      font=\scriptsize\sffamily
    ](key) {#1\strut}
  ;
}

%====================================
% Desactivación de división de palabras. 
% Uso: \nodivide o \nodivide[<n>]
\newcommand{\nodivide}[1][10000]{%
	\hyphenpenalty=#1 % Valor típico: hasta 10000
	\exhyphenpenalty=#1 % Valor típico: hasta 10000
	\sloppy
}

%====================================
% Desactivación de división de palabras. 
% Uso: \nowidowandorphan o \nowidowandorphan[<n>]
\newcommand{\nowidowandorphan}[1][10000]{%
	\clubpenalty=#1  % % Valor típico: hasta 10000
	\widowpenalty=#1 % % Valor típico: hasta 10000
}

%====================================
% Código para evitar la división de notas al pie en págs. diferentes
% Uso: \nodividenotas o \nodividenotas[<n>]
\newcommand{\nodividenotes}[1][10000]{%
	\interfootnotelinepenalty=#1 % Valor típico: hasta 10000
}


%====================================
% Creación de un contador nuevo para almacenar el nº de páginas actual
% OJO: Debe ir antes de \mainmatter (antes de que se reinicie en cnt page)
\newcommand{\savepagecnt}{%
	\newcounter{totpages}
	\setcounter{totpages}{\value{page}}
	\addtocounter{totpages}{1}
}

%====================================
% Continuación de la paginación desde el valor almacenado en \totpages
\newcommand{\contpagination}{%
	\setcounter{page}{\value{totpages}}
}

%====================================
% Limpia las cabeceras de la primera página de capítulo
\newcommand{\cleanhdfirst}{%
	\fancypagestyle{plain}{%
		\fancyhf{}%
		\renewcommand{\headrulewidth}{0pt}
		\renewcommand{\footrulewidth}{0pt}
	}
}

%====================================
% Añadido de entorno abstract en clase book
% donde no está definido por defecto
\newenvironment{abstract}%
{\cleardoublepage\null \vfill\begin{center}%
\bfseries \sffamily \abstractname \end{center}}%
{\vfill \null}


%====================================
% Comando para incluir una nota de aviso al margen.
% Si el margen no es suficientemente amplio puede generar un Overfull \hbox
\newcommand{\ojo}[1]{%
	\marginpar{\footnotesize\raggedright\ding{42} #1}}
\newcommand{\uju}{%
	\renewcommand{\ojo}[1]{}}











%%%%%%%%%%%%%%
% Datos del documento
% Estos valores son los que se emplean en el documento por lo que no son
% traducidos. Cuando algún campo puede tener varias líneas aparecen dos
% campos señalados como <campo>Primera y <campo>Segunda.
% Si no se desea emplear un campo se comenta.
%====================================
% EDITAR: Datos del documento
\tituloPrimera{<Primera línea título>}% 1ª Línea
\tituloSegunda{<Segunda línea título>}% 2ª Línea  % Comentar si no existe
\titulo{<Título corto>}% Título corto
\autor{<autor (nombre apellidos)>}
\email{<autor>@uclm.es}
\director{<director (nombre apellidos)>}
\codirector{<codirector (nombre apellidos)>} % Comentar si no existe
\instEdu{UNIVERSIDAD DE CASTILLA-LA MANCHA}
\escudo{escudoInf} % Inclusión del escudo de la institución
%\escudo{escudoUCLM} 
\centroEdu{ESCUELA SUPERIOR DE INFORMÁTICA}
\deptoEduPrimera{<Primera línea Depto. Director>}% 1ª Línea
\deptoEduSegunda{<Segunda línea Depto. Director>}% 2ª Línea  % Comentar si no existe
\titulacion{GRADO EN INGENIERÍA INFORMÁTICA}
\especialidad{<Tecnología Específica>} % Tecnología específica, ...
\tipoDoc{TRABAJO FIN DE GRADO} % Proyecto Fin de Carrera, Trabajo Fin de Grado, Tesis Doctoral, ...
% Si las fechas se desean en inglés hay que ponerlas explícitamente.
\fechaDef{<mes, año>} % Fecha de defensa
\mesDef{<mes>}        % Mes de defensa
\yearDef{<año>}        % Año de defensa
\lugarDef{<Ciudad>}% Lugar de defensa
%%%%%%%%%%%%%%
\begin{document}
% --- Bloque inicial del documento ---
\ifspanish
\selectlanguage{spanish}% Emplea idioma español
\else
\selectlanguage{english}% Emplea idioma inglés
\fi 
\frontmatter
% Cambia la numeración de páginas a números romanos y las secciones no están numeradas aunque si aparecen en el índice de contenidos.
\pagestyle{empty}  % Páginas sin cabecera ni pies

%%%%%%%%%%%%%%
\portadaTFG

\portadillaTFG
%%%%%%%%%%%%%%

% Créditos
% EDITAR: Licencia (si se desea modificar).
\creditos{Este documento se distribuye con licencia Creative Commons Atribución Compartir Igual 4.0. El texto completo de la licencia puede obtenerse en \url{https://creativecommons.org/licenses/by-sa/4.0/}. La copia y distribución de esta obra está permitida en todo el mundo, sin regalías y por cualquier medio, siempre que esta nota sea preservada. Se concede permiso para copiar y distribuir traducciones de este libro desde el español original a otro idioma, siempre que la traducción sea aprobada por el autor del libro y tanto el aviso de copyright como esta nota de permiso, sean preservados en todas las copias.}{cclicense}

%%%%%%%%%%%%%%
\tribunalTFG


%%%%%%%%%%%%%%
% Dedicatoria (opcional, 1 pág.)
% Aunque opcional, no se debería perder la oportunidad de poder dedicar el trabajo a alguien MUY especial.
% EDITAR: Dedicatoria
\dedicado{Alguien MUY especial ... \\
(para siempre)} % Editar para dedicar a alguien

% Ajustes para texto normal
% Estos tres comandos admiten un parámetro entre [corchetes] que se puede ir elevando desde 5000 (defecto) hasta 10000. Cuanto mayor es el valor más forzado está LaTeX
%\nodivide[10000]      % Evita la división de palabras 
%\nowidowandorphan % Evita viudas y huérfanas
%\nodividenotes %Evita división entre páginas de las notas
\pagestyle{plain} % Páginas sólo con numeración inferior en pie

%%%%%%%%%%%%%%
% Resumen (Si se desea cambiar el orden en que se imprimen, cambiar de forma manual)
\selectlanguage{spanish}
\begin{abstract}
	% EDITAR: Resumen
	(... versión del resumen en español ...)
El resumen debe ocupar como máximo una página y en dicho espacio proporcionará información crucial sobre el \emph{`qué'} (problemática que trata de resolver el TFG), el \emph{`cómo'} (metodología para llegar a los resultados) y los objetivos alcanzados.
\end{abstract}

% Abstract
\selectlanguage{english}
\begin{abstract}
	% EDITAR: Abstract
	(... english version of the abstract ...)
\end{abstract}


% Selección de idioma para el resto del documento
\ifspanish
\selectlanguage{spanish}% Para el resto del documento el idioma es español.
\else
\selectlanguage{english}
\fi

% Agradecimientos (1 pág.)
\chapter*{Agradecimientos} % Opción con * para que no aparezca en TOC ni numerada

%OJO: Editar
 Aunque es un apartado opcional, haremos bueno el refrán \emph{<<Es de bien nacidos, ser agradecidos>>} si empleamos este espacio es un medio para agradecer a todos los que, de un modo u otro, han hecho posible que el TFG <<llegue a buen puerto>>. Esta sección es ideal para agradecer a familiares, directores, profesores, compañeros, amigos, etc. 
 
 Estos agradecimientos pueden ser tan personales como se desee e incluir anécdotas y chascarrillos, pero nunca deberían ocupar más de una página.

\makeatletter		
\begin{flushright}
	\textit{\@autor}
\end{flushright}
\makeatother % Agradecimientos etc.




%%%%%%%%%%%%%%
% Índices
\pagestyle{fancy} % Estilo de página ajustado por fancyhdr
% OJO: Si se cambian los nombres que otorga babel para el idioma en uso, los nombres vuelven a su estado cada vez que se vuelva a seleccionar el idioma. Por eso es conveniente que el cambio de nombre aparezca justo antes de su empleo.

% Índice General
\tableofcontents  % Índice general

% Índice de figuras
%\addcontentsline{toc}{chapter}{\listfigurename} % Añade la lista de figuras al TOC.
\listoffigures    % Índice de figuras (opcional)

% Índice de tablas
\ifspanish
\renewcommand{\tablename}{Tabla} % Se sustituye 'Cuadro' por 'Tabla'
\renewcommand{\listtablename}{Índice de tablas}
\fi
%\addcontentsline{toc}{chapter}{\listtablename} % Añade la lista de tablas al TOC.
\listoftables % Índice de tablas (opcional)

% Índice de listados
% Personalización de los títulos de los listados y del índice de listados
\ifspanish
\renewcommand{\lstlistingname}{Listado} % Nombre que figura en el pie
\renewcommand{\lstlistlistingname}{Índice de listados} % Nombre asignado
\fi
%\addcontentsline{toc}{chapter}{\lstlistlistingname} % Añade la lista de listados al TOC.
\lstlistoflistings % Índice de listados creados con listings (opcional)

% Índice de algoritmos
% Personalización de los títulos de los algoritmos y del índice de algoritmos
\ifspanish
\SetAlgorithmName{Algoritmo}{Alg}{Índice de algoritmos}
\fi
%\addcontentsline{toc}{chapter}{Índice de algoritmos} % Añade la lista de algoritmos al TOC.
\listofalgorithms % Índice de algoritmos creados con algortihm2e




%%%%%%%%%%%%%%
% OPCIONAL: No es preciso incluirlo
% Lista de acrónimos o nomenclatura (Opcional)
% En libros o tesis es usual encontrarlos. En el TFG es inusual.
% Debe estar ordenado por orden alfabético.
% Me parece más interesante poner el índice alfabético final que es mucho más útil.
%\include{Nomeclatura} % (Opcional) Acrónimos y nomenclatura usada

% Prólogo (opcional)
% En libros o tesis es más usual encontrarlos. En el TFG es inusual.
%% Prólogo (opcional, 1 ó 2 pags.)
\chapter{Prólogo}
%[OJO:] Editar
%En un libro esta sección explica las motivaciones del libro y contiene habitualmente una declaración de intenciones del autor.

Este documento surgió para servir como plantilla de ejemplo en un curso, dirigido a universitarios, para la preparación de proyectos fin de carrera, tesis y otros documentos académicos con \LaTeX{}.

\begin{flushright}
	\makeatletter
	\@lugarDef, \@mesDef{} de \@yearDef\\[1cm]
	\@autor{}\\
	\@email
	\makeatother
\end{flushright} % Prólogo (opcional).




%%%%%%%%%%%%%%
% Capítulos del documento
% Creación de un contador nuevo para almacenar el nº de páginas actual
% OJO: Debe ir antes de \mainmatter (antes de que se reinicie el cnt page)
\savepagecnt
\mainmatter
% Justo antes del primer capítulo del libro. Activa la numeración con números arábigos y reinicia el contador de páginas.

% Se incluye un fichero para cada capítulo. Se emplea la instrucción \include porque en los libros lo más habitual es que cada cápítulo comience en una nueva página.

% OJO: No cambiar de ubicación. Redefinición de estilo plain para que la 1ª página del capítulo no tenga ni cabecera ni pie
\cleanhdfirst

% Reajuste del número de página consecutivo para no reiniciar paginación en Cáp. 1
%\contpagination % Comentar si se desea reiniciar paginación

% Con el comando \include se garantiza que el capítulo comienza en nueva página.

% Aquí irían todos los capítulos del documento
\chapter{Introducción}
\label{cap:Introduccion}

Introducción al tema, entorno y justificación de la importancia del trabajo abordado.

Estructura del resto del documento.

% ------------------------------------------------------------------------------
% Ejemplos para la plantilla
% ------------------------------------------------------------------------------
\section{Ejemplos de listas}
\label{sec:ejListas}
\index{ejemplos} % Véase cómo se incluyen entradas en el índice alfabético
A continuación se van a añadir algunos ejemplos que pueden emplearse al redactar la memoria.

\index{ejemplos!listas} % Para el índice
\noindent Ejemplo de lista con \emph{bullet} especial. 
% Ejemplo: Lista con bullets especiales
% ============
\begin{itemize}
	\item[*] peras
	\item manzanas
	\item[\ding{170}] naranjas
\end{itemize}

\noindent Ejemplo de lista compacta (también se puede emplear el entorno para enumeraciones \emph{compactenum})
% Ejemplo: Lista con balas
% ============
\begin{compactitem}
	\item peras
	\item manzanas
	\item naranjas
\end{compactitem}


\noindent Ejemplo de lista en varias columnas.
% Ejemplo: Listas en varias columnas
% ============
\begin{multicols}{2} % El parámetro es el número de columnas de la lista
	\begin{compactenum}
		\item peras
		\item manzanas
		\item naranjas
		\item patatas
		\item calabazas
		\item fresas
	\end{compactenum}
\end{multicols}


\newpage


\section{Ejemplos de tablas}
\label{sec:ejTablas}
\index{ejemplos!tablas}
A continuación se incluyen algunos ejemplos de tablas hechas con \LaTeX{} y paquetes dedicados.

% Ejemplo: Tabla con macro \cline
% ==========
\begin{table}[htb]%
	\centering
	\caption{Ejemplo de uso de la macro \texttt{cline}}
	\label{tab:cline}
	\begin{tabular}[t]{|r|l|}
		\hline
		7C0 & hexadecimal \\[1cm] % Ejemplo de separación fijada entre líneas
		3700 & octal \\ \cline{2-2}
		11111000000 & binario \\
		\hline \hline
		1984 & decimal \\
		\hline
	\end{tabular}
\end{table}


\noindent Ejemplo de tabla en la que se controla el ancho de la celda.

% Ejemplo: Ejemplo de tabla con control de la anchura de celda.
% ==========
\begin{table}[htb]%
	\centering
	\caption{Ejemplo de tabla con especificación de anchura de columna}
	\label{tab:anchura}
	\begin{tabular}{ | l | l | l | p{5cm} |}
		\hline
		Día & Temp Mín (\textdegree C) & Temp Máx (\textdegree C) & Previsión \\ \hline
		Lunes & 11 & 22 & Día claro y muy soleado. Sin embargo, la brisa de la tarde puede hacer que las temperaturas desciendan \\ \hline
		Martes & 9 & 19 & Nuboso con chubascos en muchas regiones. En Cataluña claro con posibilidad de bancos nubosos al norte de la región \\ \hline
		Miércoles & 10 & 21 & La lluvía continuará por la mañana pero las condiciones climáticas mejorarán considerablemente por la tarde\\
		\hline
	\end{tabular}
\end{table}

\clearpage


\section{Ejemplos de figuras}
\label{sec:ejFiguras}
\index{ejemplos!figuras}

En esta sección se añaden ejemplos de muestra para la inclusión de figuras simples y subfiguras.

% Ejemplo: Ejemplo de inclusión de figura
% ============
\begin{figure}[htb]
	\centering
	\includegraphics[width=0.4\linewidth]{escudoInf}
	\caption[Ejemplo de figura]{Figura vectorial del escudo de la ESI}
	\label{fig:ejFigure}
\end{figure}


\noindent Ejemplo de figuras compuestas por subfiguras.

% Ejemplo: Ejemplo de inclusión de subfiguras
% ============
\begin{figure}[htb]
	\centering
	\subfigure[Gráfico vectorial PDF]{
		\includegraphics[width=0.4\linewidth]{escudoInf}
		\label{fig:escudoColor}
	} 
	\subfigure[Gráfico png]{
		\includegraphics[width=0.4\linewidth]{escudoInfBW}
		\label{fig:escudoBW}
	}
	\caption[Ejemplo de subfiguras]{Ejemplo de inclusión de subfiguras en un mismo entorno}
	\label{fig:ejSubfigures}
\end{figure}


\clearpage


\section{Ejemplos de listados}
\label{sec:ejListados}
\index{ejemplos!listados}

Ejemplos más representativos de inclusión de porciones de código fuente.

% Ejemplo: Listado Java
% ============
\begin{lstlisting}[language=Java,float=ht,caption={[Código fuente en Java]Ejemplo de código fuente en lenguaje Java},label=lst:java]
// @author www.javadb.com
public class Main {    
// Este método convierte un String a
// un vector de bytes

public void convertStringToByteArray() {

String stringToConvert = "This String is 15";      
byte[] theByteArray = stringToConvert.getBytes();        
System.out.println(theByteArray.length);        
}

// argumentos de línea de comandos 
public static void main(String[] args) {
new Main().convertStringToByteArray();
}
}
\end{lstlisting}



\noindent Otro ejemplo.

\begin{lstlisting}[style=C,float=ht,caption={Ejemplo de código C},label=lst:codC]
// Este código se ha incluido tal cual está 
// en el fichero \LaTeX{}
#include <stdio.h>
int main(int argc, char* argv[]) {
puts("¡Hola mundo!");
}
\end{lstlisting}


\noindent Ejemplo de entrada por consola.

\begin{lstlisting}[style=consola, numbers=none]
$ gcc -o Hola HolaMundo.c
\end{lstlisting}
\chapter{Objetivo}
\label{cap:Objetivo}

Introduce y motiva la problemática (i.e.\emph{\ ¿cuál es el problema que se plantea y porqué es interesante su resolución?})

Debe concretar y exponer detalladamente el problema a resolver, el entorno de trabajo, la situación y qué se pretende obtener. También puede contemplar las limitaciones y condicionantes a considerar para la resolución del problema (lenguaje de construcción, equipo físico, equipo lógico de base o de apoyo, etc.). Si se considera necesario, esta sección puede titularse \emph{Objetivos del TFG e hipótesis de trabajo}. En este caso, se añadirán las hipótesis de trabajo que el alumno pretende demostrar con su TFG.

Una de las tareas más complicadas al proponer un TFG es plantear su \textsf{Objetivo}. La dificultad deriva de la falta de consenso respecto de lo que se entiende por \emph{objetivo} de un trabajo de esta naturaleza. En primer lugar se debe distinguir entre dos tipos de objetivo:

\begin{enumerate}
	\item La \emph{finalidad específica} del TFG que se plantea para resolver una problemática concreta aplicando los métodos y herramientas adquiridos durante la formación académica. Por ejemplo, \emph{<<Desarrollo de una aplicación software para gestionar reservas hoteleras \emph{on-line}>>}.
	
	\item El \emph{propósito académico} que la realización de un TFG tiene en la formación de un graduado. Por ejemplo, la \emph{adquisición de competencias específicas de la especialización} cursada.
\end{enumerate}

En el ámbito de la memoria del TFG se tiene que definir el primer tipo de objetivo, mientras que el segundo tipo de objetivo es el que se añade al elaborar la propuesta de un TFG presentada ante un comité para su aprobación. Este segundo tipo de objetivo no debe incluirse en el apartado correspondiente de la memoria y en todo caso puede valorarse su satisfacción en la sección de resultados y conclusiones.

Un objetivo bien planteado para el TFG debe estar determinado en términos del \emph{<<producto final>>} esperado que resuelve un problema específico. Es por tanto un sustantivo que debería ser \emph{concreto} y \emph{medible}. El \textsf{Objetivo} planteado puede pertenecer una de las categorías que se indica a continuación:
\begin{itemize}
	\item \emph{Diseño y desarrollo de <<artefactos>>} (habitual en las ingenierías),
	\item \emph{Estudio} que ofrece información novedosa sobre un tema (usual en las ramas de ciencias y humanidades), y
	\item \emph{Validación de una hipótesis} de partida (propio de los trabajos científicos y menos habitual en el caso de los TFG).
\end{itemize}

Estas categorías no son excluyentes, de modo que es posible plantear un trabajo cuyo objetivo sea el diseño y desarrollo de un <<artefacto>> y éste implique un estudio previo o la validación de alguna hipótesis para guiar el proceso. En este caso y cuando el objetivo sea lo suficientemente amplio puede ser conveniente su descomposición en elementos más simples hablando de \emph{subobjetivos}. Por ejemplo, un programa informático puede descomponerse en módulos o requerir un estudio previo para plantear un nuevo algoritmo que será preciso validar. 

La descomposición de un objetivo principal en subobjetivos u objetivos secundarios debería ser natural (no forzada), bien justificada y sólo pertinente en los TFG de gran amplitud.

Junto con la definición del objetivo del TFG se puede especificar los \emph{requisitos} que debe satisfacer la solución aportada. Estos requisitos especifican \emph{características} que debe poseer la solución y \emph{restricciones} que acotan su alcance. En el caso de TFG cuyo objetivo es el desarrollo de un <<artefacto>> los requisitos pueden ser \emph{funcionales} y \emph{no funcionales}.

Al redactar el objetivo de un TFG se debe evitar confundir los medios con el fin. Así es habitual encontrarse con objetivos definidos en términos de las \emph{acciones} (verbos) o \emph{tareas} que será preciso realizar para llegar al verdadero objetivo. Sin embargo, a la hora de planificar el desarrollo del trabajo si es apropiado descomponer todo el trabajo en \emph{hitos} y estos en \emph{tareas} para facilitar dicha \emph{planificación}.

%El motivo de tal confusión deriva de <<mezclar>> la definición de la finalidad específica del TFG con el propósito académico del mismo.

La categoría del objetivo planteado justifica modificaciones en la organización genérica de la memoria del TFG. Así en el caso de estudios y validación de hipótesis el apartado de resultados y conclusiones debería incluir los resultados de experimentación y los comentarios de cómo dichos resultados validan o refutan la hipótesis planteada.


\chapter{Antecedentes}
\label{cap:Antecedentes}

En esta parte se deben mostrar los conocimientos obtenidos en la búsqueda bibliográfica y no ideas personales del autor; como mucho podrá aportar un comentario de algunas pocas ideas extraídas de las fuentes en que se ha basado. Se articulará esta parte en diversas secciones, que permitan la exposición estructurada y didáctica de los conocimientos de la investigación bibliográfica.
\chapter{Metodología}
\label{cap:Metodologia}

En este capítulo se debe detallar las metodologías empleadas para planificación y desarrollo del trabajo,
así como explicar de modo claro y conciso cómo se han aplicado dichas metodologías.


Obtención de la BBDD.

Ante la nula información almacenada de una forma estructurada se ha tenido que buscar
 diversas maneras de obtener y almacenar la información, estructurarla y tratarla para
 poder sacar conocimiento de ella.

En la página de la federación internacional de esgrima (FIE) se almacenan los resultados
 de las competiciones de los últimos tiempos. De ahí se puede obtener los resultados
 de cada competición tanto como el ranking general, pasando por la fase de poules
 y acabando con los tablones eliminatorios de estos mismos. Viendo toda esta información
 almacenada se decide extraer y almacenar dicha información. Para la extracción se
 utilizarán técnicas de scrapping web mediante la cual se descarga la página y se
 obtiene su contenido para tratarlo. En este caso navegamos por dicha página y una
 vez obtenida la información la guardamos en nuestra BBDD. Puesto que no está del
 todo estructurada tendremos que ir almacenando dicha información en diferentes BBDD
 y después juntarlas.

Descripción de la BBDD.

La primera BBDD que generaremos será aquella en la que guardaremos la información
 básica de los asaltos. Para ello guardaremos el ID de la competición, el número
 de tablón en el que se efectuó el asalto, ID del primer tirador, ID del segundo
 tirador, tocados obtenidos por el segundo primer tirador y tocados obtenidos por
 el segundo tirador.

\begin{table}[htb]%
  \centering
  \caption{Estructura BBDD asaltos inicial}
  \label{tab:anchura}
  \begin{tabular}{ | l | l | l | l | l | l | }
    \hline
    Nombre de Campo & Tipo de campo & Ejemplo \\ \hline
    CompetitionID & Texto & 2019-64 \\ \hline
    Tableu & Entero & 32 \\ \hline
    Competitor1 & Texto & /fencers/Anna-KOROLEVA-40351/ \\ \hline
    Competitor2 & Texto & /fencers/Kira-KESZEI-49034/ \\ \hline
    ResultCompetitor1 & Texto & V/15 \\ \hline
    ResultCompetitor2 & Texto & D/3 \\
    \hline
  \end{tabular}
\end{table}

El siguiente paso que tendremos que dar será obtener la información de los tiradores
 para ello ser visitará la página correspondiente. Un ejemplo de página que almacena
 la información de un tirador sería el siguiente http://fie.org/es/fencers/Mario-PERSU-31870
 como se puede observar el final de la URL es el mismo que el identificador almacenado en la
 anterior tabla. De modo que explorando la anterior BBDD podremos visitar las páginas de cada
 tirador almacenado y de ese modo generar una nueva BBDD con toda la información de cada uno de ellos.
 En dicha BBDD tendremos su identificador, edad, ranking (si lo tuvieran), nacionalidad, mano usada
 y arma.

\begin{table}[htb]%
  \centering
  \caption{Estructura BBDD tiradores}
  \label{tab:anchura}
  \begin{tabular}{ | l | l | l | l | l | l | }
    \hline
    Nombre de Campo & Requerido & Tipo de campo & Ejemplo \\ \hline
    competitorID & Si & Texto & ADRIANA-MILANO-36467 \\ \hline
    Age & Si & Entero & 21 \\ \hline
    FieRanking & No & Entero & 300 \\ \hline
    HandNess & Si & Texto & Right \\ \hline
    Weapon & Si & Texto & Epée \\
    \hline
  \end{tabular}
\end{table}

Una vez obtenidos todos los datos referentes a los tiradores tendremos que cruzar
 las dos tablas mencionadas anteriormente de modo que tengamos toda la información
 en una sola BBDD. Esta última BBDD tendrá la información de la primera, sustituyendo
 las columnas de ID de cada tirador por su registro correspondiente en la anterior tabla.

De modo que la estructura será la siguiente
\begin{table}[htb]%
  \centering
  \caption{Estructura BBDD asaltos final}
  \label{tab:anchura}
  \begin{tabular}{ | l | l | l | l | l | l | }
    \hline
    Nombre de Campo & Tipo de campo & Ejemplo \\ \hline
    ComptetitionID & Texto & 2019-176 \\ \hline
    TABLEU & Entero & 32 \\ \hline
    C1\_ID & Texto & Sergey-KHODOS-13869 \\ \hline
    C1\_RANKING & Entero & 22 \\ \hline
    C1\_NATIONALITY & Texto & RUSSIA \\ \hline
    C1\_HANDNESS & Texto & Right \\ \hline
    C1\_WEAPON & Texto & Epée \\ \hline
    C2\_ID & Texto & Laurin-EGGENSCHWILER-5966 \\ \hline
    C2\_RANKING & Entero & 34 \\ \hline
    C2\_NATIONALITY & Texto & SWITZERLAND \\ \hline
    C2\_HANDNESS & Texto & Right \\ \hline
    C2\_WEAPON & Texto & Epée \\ \hline
    C2\_ID & Texto & Laurin-EGGENSCHWILER-5966 \\ \hline
    C1\_RESULT & Texto & V/15 \\ \hline
    C2\_RESULT & Texto & D/3 \\
    \hline
  \end{tabular}
\end{table}

Modo de empleo.

El modo de empleo de esta BBDD será para el entrenamiento de una red neuronal
 la cual servirá para apoyar al sistema experto. En total hay unos 28.000 registros,
 de los cuales se emplearán entorno al 60 por ciento para entrenar al modelo y un 40
 por ciento para comprobar que el entrenamiento ha sido satisfactorio.

\chapter{Resultados}
\label{cap:Resultados}

En esta sección se describirá la aplicación del método de trabajo presentado en el capítulo \ref{cap:Metodologia} en este caso concreto, mostrando los elementos (modelos, diagramas, especificaciones, etc.) más importantes. Este apartado debe explicar cómo la metodología satisface los objetivos y requisitos planteados.
\chapter{Conclusiones}
\label{cap:Conclusiones}

Breve resumen de lo más destacable del TFG con la solución propuesta y posibles mejoras, ampliaciones o trabajos relacionados que quedan por hacer y que tienen interés para el tema tratado.





%%%%%%%%%%%%%%
% Anexos
\appendix

% Tras este punto los capítulos se numeran con letras.
\ifspanish
\renewcommand{\appendixname}{Anexo}
\fi
\chapter{Anexo A}
\label{cap:Test Slagel}


pendiente escribir


 % Apéndice A (opcionales)
%
%............... % Aquí todos los apéndices necesarios


%%%%%%%%%%%%%%
% Bibliografía
\backmatter

\nocite{*} % Comando por si queremos incluir todas las entradas de la bibliografía aunque no estén citadas en el texto pral.(NO ESTÁ PERMITIDO POR LA NORMATIVA DEL ALGUNOS CENTROS (p.ej. ESI)-> Todas las referencias deben estar citadas en el texto)

\ifspanish
\renewcommand{\bibname}{Bibliografía}
\fi
\addcontentsline{toc}{chapter}{\bibname} % Añade la bibliografía al Índice de contenidos.
\printbibliography



%%%%%%%%%%%%%%
% Índice alfabético (opcional)
% OPCIONAL: (Consejo) Incluir los comandos mientras se escribe cada capítulo ya que hacerlo al final resulta tedioso.
\cleardoublepage
\ifspanish
%\renewcommand{\indexname}{Índice alfabético} %Cambio de nombre de sección
\fi
\addcontentsline{toc}{chapter}{\indexname} % Añade al Índice de contenidos.
\printindex  % Facilitado por makeidx (opcional, si no se usa no se imprime)
\end{document}